\section{Introduction}
\label{sec:introduction}


Neural scaling laws have ushered in new paradigms for interpreting deep learning dynamics. By establishing that performance gains follow predictable behavior relative to model size, dataset scale, and compute, neural scaling laws have pushed stochastic optimization methods toward increasingly extended training horizons under tight efficiency constraints~\cite{hestness2017deep, kaplan2020scaling, henighan2020scaling, hoffmann2022chinchilla}. In this regime, optimization stability becomes an important practical requirement, since unstable training may result in slow convergence and critical loss spikes~\cite{bai2026spikes}. In practice, stability is often maintained through a set of heuristics, such as learning-rate decay and gradient clipping~\cite{loshchilov2016sgdr,pascanu2013difficulty,bai2026spikes}. Yet, recurrent loss spikes can coexist with an overall decrease in training loss over longer timescales.~\cite{cohen2021edge,damian2023self}. Reducing the learning rate or limiting gradient magnitude can suppress loss spikes~\cite{smith2017cyclical,pascanu2013difficulty}, but smoother loss trajectories do not necessarily imply faster optimization. The central question is therefore: which changes should a learning intervention target to improve training?

Training becomes unstable when updates are too large relative to sharp directions of the loss landscape, so that they overshoot rather than converge towards a minimum~\cite{polyak1964methods, bertsekas1997nonlinear}. However, the Edge-of-Stability literature has shown that deep networks often continue to make progress near, or even beyond, the nominal sharpness threshold, suggesting that curvature alone is not a sufficient explanation~\cite{cohen2021edge, arora2022understanding, damian2023self}. An adaptive optimizer such as Adam~\cite{kingma2014adam} adds additional complexity through memory and coordinate scaling, and more recently it has been shown that its adaptive preconditioner can produce loss spikes even in a quadratic loss function~\cite{bai2026spikes}.
%Stochastic extensions, such as the break-even point framework, argue that instability depends jointly on sharpness and curvature variance, particularly during an early phase in which learning rate and batch size strongly shape the subsequent trajectory~\cite{jastrzebski2020break}. 

These studies are limited by their reliance on loss-landscape properties alone to explain optimization and instability. The loss landscape is highly complex, problem-specific~\cite{ly2025multifractal,yang2021taxonomizing,ma2022multiscale} and closely tied to network architecture~\cite{li2018visualizing}. This motivates examining how parameter updates affect representations within individual layers, beyond their aggregate effect on the loss. More recently, Central Flows has explicitly highlighted the need to connect loss dynamics to the roles of individual layers~\cite{cohen2025central}. This work addresses this gap. We argue that the network should not be treated solely as a black box, and we connect learning instability and changes in hidden representations directly by measuring how updates alter internal responses.

In summary, we make four main contributions:
\begin{enumerate}
\item \textbf{Theoretical and empirical characterization of constructive and destructive regimes.} Starting from the SGD update equations, we express net parameter displacement through gradient magnitudes and alignment throughout the entire learning trajectory. We characterize constructive regimes by persistent reinforcement between consecutive updates (\(\rho_t>0\)) and destructive regimes by persistent reversal (\(\rho_t<0\)), in which updates partially undo preceding learned structure. Across MLP, ResNet, and Transformer, destructive regimes exhibit back-and-forth motion while training loss continues to decrease. We quantify the fraction of total distance traveled that remains as net parameter displacement through \(R_T\), which falls to approximately \(2.6\)--\(5.4\%\). We further characterize the effect of mini-batch sampling: smaller batches attenuate gradient alignment in magnitude, dampening the oscillatory motion. Finally, we connect gradient reversal and amplification to loss increases: for quadratic full-batch GD, a completed step increases loss precisely when \(\rho_t<-\|g_{t-1}\|/\|g_t\|\). We show that this is equivalent to the directional-curvature loss criterion, but described through consecutive-gradient metrics. This condition is general across datasets and architectures.

\item \textbf{Recovering directional curvature from gradients diagnostics.} Bai et al.~\cite{bai2026spikes} highlight the computational cost of computing Hessian eigenvalues as a limitation for practical loss-spike analysis. Building on classical secant geometry~\cite{barzilai1988two,tan2016bbsgd}, we show that consecutive gradients and the applied parameter update recover directional curvature averaged over the completed step, exactly for smooth, fixed objectives and independently of the optimizer. This diagnostic reuses training gradients without additional backward passes or Hessian--vector products. We further investigate whether \(\rho_t<-\|g_{t-1}\|/\|g_t\|\) provides advance warning of loss increases under Adam. Across MLP, ResNet, and Transformer, it achieves higher accuracy and warning lift than constant predictions, the preceding loss change, and rules based on alignment or norm growth alone when predicting a next-step loss increase.

\item \textbf{Functional characterization and tracing of constructive and destructive regimes.} We extend our characterization from gradient interactions to changes in hidden representations and network outputs, addressing the need to connect optimization dynamics to individual layers, highlighted by Central Flows~\cite{cohen2025central}. Empirically, we identify approximately low-dimensional functional oscillations across multiple updates. Furthermore, net output displacement retains only approximately \(1.7\)--\(3.5\%\) of the total distance traveled. These oscillations exhibit shared temporal structure across network sites, but their distribution through the network differs across architectures: reversal is homogeneous across all MLP sites, whereas earlier ResNet and Transformer representations can remain constructive while their outputs reverse. Together, these findings establish a functional counterpart to the gradient-based characterization of destructive regimes across multiple datasets and architectures.

\item \textbf{Testing representation displacement as an intervention target.} Kosson et al.~\cite{kosson2024warmup} identify representation change as a promising measure of update size, but leave its effective control as an open challenge. Motivated by our functional analysis, we introduce a distributed displacement cap that scales down Adam's proposed updates according to the direct response changes they would produce. We compare this intervention with unchanged Adam, learning-rate reduction, second-moment changes, and gradient clipping. Across all models, every tested cap quantile reduces mean training loss and accumulated upward loss spikes relative to unchanged Adam. 

\end{enumerate}

The remainder of the paper is organized as follows. Section~\ref{sec:background} positions these contributions in related work. Sections~\ref{sec:displacement}--\ref{sec:gd-experiments} develop and test the GD/SGD framework; Sections~\ref{sec:adam}--\ref{sec:vc-interventions} connect Adam diagnostics to functional measurements and interventions. Section~\ref{sec:conclusion} discusses implications, limitations and future work.